\section{Related Work}
\label{sec:related-work}

GAWorld draws on four bodies of work that answer different parts of the
artificial-society problem.  Its contribution is a system integration and an
audited set of research surfaces; we do not claim that the underlying ideas of
memory, tool-using agents, social evaluation, or agent-based validation are
unique to GAWorld.

\subsection{Generative agents and simulated individuals}

Generative Agents demonstrated that an LLM-driven architecture combining an
experience stream, retrieval, reflection, and planning can sustain socially
legible behavior in a sandbox community \citep{park2023generative}.  That work
provides a central precedent for persistent social characters and for studying
emergent interaction through an inspectable environment.  Later work grounds
agents in extensive self-report data and evaluates how well responses predict
held-out answers from the same individuals \citep{park2024grounded}.  Research
on ``silicon samples'' similarly explores conditioning language models to
approximate distributions of survey responses
\citep{argyle2023outofone}.  These studies motivate profile grounding and
longitudinal state, but they also distinguish two validation targets:
internally coherent simulated inhabitants and externally evaluated models of
particular people or populations.

GAWorld is oriented toward constructing the surrounding research system.  It
adds configurable urban, social, economic, work, information, and policy
mechanisms; multiple operational interfaces; and explicit experiment and
artifact workflows around a persistent-agent loop.  Its archived profiles and
traces have not undergone the held-out, person-level validation used in work
on self-report-grounded agents.  Accordingly, GAWorld adopts persistence as a
system property and treats fidelity as an open empirical claim.

\subsection{Language-agent cognition and architecture}

LLM-agent research has developed reusable patterns for coupling language
generation to action.  ReAct interleaves reasoning and environment-facing
actions \citep{yao2023react}; Reflexion stores verbal feedback to improve later
attempts \citep{shinn2023reflexion}; and Voyager combines iterative prompting,
environment feedback, and an accumulating skill library in an open-ended
embodied setting \citep{wang2023voyager}.  MemoryBank studies long-term memory
with retrieval and forgetting mechanisms \citep{zhong2024memorybank}, while
CoALA offers a conceptual organization of language agents in terms of memory,
action space, and decision procedures \citep{sumers2024coala}.  CAMEL, in
turn, examines role-conditioned communication among language agents
\citep{li2023camel}.

These systems and abstractions inform GAWorld's separation of structured
state, episodic storage, retrieval, planning, action, reflection, and learned
skills.  GAWorld differs in scope rather than by proposing a new generic
reasoning algorithm: it connects such components to a persistent multi-agent
society and to a clocked, file-backed experimental runtime.  Its twelve-stage
pipeline, event bus, and plugin registry expose replacement and observation
points, but the present migration remains mixed with direct and hook-based
paths.  Thus GAWorld contributes an implemented integration for social
simulation, not a claim of a universally superior cognitive architecture.

\subsection{Social-agent tasks and evaluation}

AgentBench evaluates LLM agents across interactive environments and emphasizes
the gap between language ability and reliable action
\citep{liu2024agentbench}.  AgentBoard adds fine-grained progress measures and
analytical visualization for multi-turn agents \citep{ma2024agentboard}.
SOTOPIA focuses specifically on social intelligence by placing agents in
interactive social scenarios and evaluating goal achievement and social
dimensions \citep{zhou2024sotopia}.  Together, these benchmarks provide
controlled tasks, comparative evaluation, and diagnostics for an agent or a
small interaction.

GAWorld addresses a complementary unit of analysis: a persistent society run
whose agents, environment, institutions, and artifacts evolve across simulated
time.  Its benchmark layer can check selected outputs, known-sign diagnostics,
placebos, and determinism when suitable artifacts exist, but it is not a
replacement for standardized social-intelligence or interactive-agent
benchmarks.  Conversely, task success alone does not validate population
composition, emergent dynamics, or policy counterfactuals.  GAWorld therefore
retains agent-level traces and run-level evidence status so that task
competence, social interaction, and society-level claims can be evaluated
separately.

\subsection{Agent-based social simulation and validation}

Agent-based modeling (ABM) treats macro patterns as outcomes of specified
micro-level rules and interactions, while generative social science emphasizes
using such models to explain how collective regularities can arise
\citep{epstein2006generative}.  The ODD protocol provides a standard vocabulary
for documenting an individual- or agent-based model's overview, design
concepts, and details \citep{grimm2006odd}.  Methodological work on empirical
ABM validation stresses that matching selected aggregate observations is not
sufficient: calibration, identification, sensitivity, multiple validation
targets, and out-of-sample evidence remain distinct concerns
\citep{windrum2007empirical,fagiolo2007critical}.

GAWorld inherits this demand for explicit mechanisms and layered evidence, but
LLM mediation introduces additional sources of variation.  Prompt context,
provider behavior, model versions, parsing fallbacks, and generated text can
change trajectories even when conventional parameters and seeds are recorded.
The system therefore couples ABM-style configuration and intervention design
to provider logs, raw agent traces, per-arm completion status, and an artifact
ledger.  Its present archive supports software-path and affordance claims, not
empirical validation of a population model.  This distinction also responds to
concerns that language models should not be treated as drop-in replacements
for human participants \citep{dillion2023replace}, particularly when generated
outputs can flatten variation within identity groups
\citep{wang2025flatten}.
